% =====================================================================
% 1. GİRİŞ (Introduction)
% Evidence base: project_plan.md §1, project_structure.md §2.4 (verified facts),
%                docs/school_assignment.md, docs/decisions.md,
%                docs/report/report_comparison_framework.md, references/INDEX.md.
% Rules: no results, no SOTA, macro-F1 headline, accuracy supporting only.
% =====================================================================
\section{Giriş}
\label{sec:introduction}

Bu projede ele alınan problem, gastrointestinal endoskopi görüntülerinin
HiperKvasir veri kümesindeki 23 sınıf üzerinde otomatik olarak
sınıflandırılmasıdır~\cite{borgli2020hyperkvasir}. Görevin temel zorluğu yalnızca
yüksek genel doğruluk elde etmek değil, sınıflar arasında oldukça dengesiz bir
dağılımın bulunduğu bir ortamda model davranışını anlamlı biçimde
değerlendirebilmektir. Bu nedenle çalışmada başarı ölçütü olarak tek başına
doğruluk (accuracy) değil, sınıf başına performansa duyarlı makro-F1 öne
çıkarılmıştır.

% --- Veri kümesinin zorluğu / dengesizlik ---
HiperKvasir etiketli alt kümesi 23 sınıf ve toplam 10.662 görüntü
içerir~\cite{borgli2020hyperkvasir}. Sınıf dağılımı ciddi biçimde dengesizdir:
en kalabalık sınıf (\textit{polyp}) yaklaşık 1028 örnek içerirken, en nadir
sınıf (\textit{hemorroids}) yalnızca 6 örnekle temsil edilmektedir.
Bu tür bir dağılımda, çoğunluk sınıflarında elde edilen yüksek doğruluk, nadir
sınıflardaki başarısızlığı gizleyebilir. Aggregate metrikler bu nedenle
yanıltıcı olabilir ve sınıf dengesizliğini açıkça dikkate alan bir değerlendirme
gerektirir.

% --- Çoklu CNN füzyonu motivasyonu ---
Farklı CNN mimarileri, aynı görüntüden farklı tümleyici temsiller öğrenme
eğilimindedir. Bu çalışmada üç önceden eğitilmiş omurga ağ --- ResNet50,
MobileNetV2 ve EfficientNetB0 --- özellik çıkarıcı olarak kullanılmış ve bu
ağlardan elde edilen özellik vektörleri birleştirilerek (özellik füzyonu) tek
bir temsil oluşturulmuştur. Buradaki temel soru, birden fazla omurga ağın
özelliklerini birleştirmenin tek bir omurga ağa kıyasla daha ayırt edici bir
temsil sağlayıp sağlamadığıdır. Gastrointestinal endoskopide çoklu-CNN özellik
füzyonu daha önce de incelenmiştir~\cite{ramamurthy2022effimix, ahmed2023hybrid};
ancak bu çalışmalar farklı protokoller kullandığından (Bölüm~\ref{sec:results})
yalnızca bağlamsal birer öncül olarak anılır, doğrudan karşılaştırma yapılmaz.

% --- Kontrollü karşılaştırma eksenleri ---
Çalışma, tek bir modelin performansını raporlamak yerine kontrollü
karşılaştırmalar üzerine kuruludur. Üç eksende karşılaştırma yapılır:
\begin{itemize}
  \item \textbf{Omurga ağ karşılaştırması:} ResNet50, MobileNetV2 ve
        EfficientNetB0'ın tek başına özellik çıkarıcı olarak performansı.
  \item \textbf{Füzyon karşılaştırması:} tek / ikili / üçlü omurga
        birleşimleri ile zorunlu iki füzyon yöntemi olan birleştirme
        (concatenation) ve ağırlıklı füzyonun (weighted fusion)
        karşılaştırılması; ek olarak gelişmiş bir kapılı füzyon (GMU) varyantı.
  \item \textbf{Transfer öğrenme karşılaştırması:} özelliklerin sabit tutulduğu
        (frozen) çıkarım ile son katman bloklarının ince ayarlandığı
        (fine-tuning) yaklaşımın karşılaştırılması.
\end{itemize}
Sınıflandırıcı tüm bu karşılaştırmalarda sabittir ve proje kararları gereği
yalnızca bir çok katmanlı algılayıcı (MLP) olarak seçilmiştir; bu nedenle
sınıflandırıcı bazlı bir karşılaştırma yapılmamıştır. Ödev metnindeki
``3 sınıflandırıcı'' ifadesinin sehven yazıldığı eğitmen tarafından teyit
edilmiştir; sınıflandırıcı MLP olarak sabitlenmiştir (PLD-06) ve karşılaştırmalar
yalnızca yukarıdaki üç eksende yapılır.

Bu üç temel eksene ek olarak, en iyi yapılandırmanın üzerinde toplamsal
(additive) ablasyonlar da değerlendirilmiştir: sınıf dengesizliğine duyarlı
focal loss, test zamanı veri artırma (TTA) ve sızıntısız seed topluluğu. Bu
ablasyonların hiçbiri, \%95 güven aralığı düzeyinde Week 3 çapraz-doğrulama
şampiyonunu (CE tabanlı en iyi yapılandırma) istatistiksel olarak geçememiştir;
final model olarak bu şampiyonun TTA'lı sürümü dondurulmuştur. Sayısal sonuçlar
ve dürüst değerlendirme Bölüm~\ref{sec:results} ve Bölüm~\ref{sec:discussion}'de
verilir.

% --- Metrik tercihi ve protokol ---
Değerlendirmede ana ölçüt olarak makro-F1 kullanılmış, doğruluk ise yalnızca
destekleyici bir ölçüt olarak raporlanmıştır. Makro-F1, her sınıfa eşit ağırlık
vererek nadir sınıflardaki başarısızlığı görünür kıldığı için dengesiz çok
sınıflı bir görevde daha bilgilendiricidir. Tüm karşılaştırmalar, sızıntıya
karşı denetlenmiş resmi HiperKvasir 5-katlı (5-fold) protokolü altında
yürütülmüş; böylece yapılandırmalar arasındaki karşılaştırmalar aynı bölme
mantığı üzerinde, metodolojik olarak tutarlı biçimde yapılmıştır.

% NOTE: Bu bölümün son paragrafı (projenin ana katkısına dair öğrencinin kendi
%       cümleleri) AGENT_REPORT_RULES.md §9 gereği öğrenci tarafından elle
%       gözden geçirilecektir.
